\documentclass[11pt]{article}

\usepackage{amsmath}
\usepackage{amssymb}

\begin{document}

\centerline{\Large\bf Math Reference}\par
\bigskip

\section{Inline Math}

Use dollar signs for inline math: $a^2 + b^2 = c^2$.
Subscripts and superscripts: $x_i$, $x^2$, $x_i^2$.
Common constants: $\pi$, $e$, $\infty$.

\section{Display Math}

Use \texttt{equation} for numbered display math:
\begin{equation}
    E = mc^2
\end{equation}

Use \texttt{equation*} (or \texttt{\$\$...\$\$}) to suppress numbering:
\begin{equation*}
    \int_0^\infty e^{-x^2}\, dx = \frac{\sqrt{\pi}}{2}
\end{equation*}

\section{Fractions and Roots}

\begin{equation*}
    \frac{a}{b} \qquad \frac{x^2 + 1}{2x - 3} \qquad \sqrt{x} \qquad \sqrt[3]{x}
\end{equation*}

\section{Sums, Products, Integrals}

\begin{equation*}
    \sum_{i=1}^{n} i = \frac{n(n+1)}{2}
    \qquad
    \prod_{i=1}^{n} i = n!
    \qquad
    \int_a^b f(x)\, dx
\end{equation*}

\section{Limits}

\begin{equation*}
    \lim_{x \to 0} \frac{\sin x}{x} = 1
    \qquad
    \lim_{n \to \infty} \left(1 + \frac{1}{n}\right)^n = e
\end{equation*}

\section{Common Symbols}

Relational: $\leq$, $\geq$, $\neq$, $\approx$, $\equiv$, $\propto$ \\
Set notation: $\in$, $\notin$, $\subset$, $\subseteq$, $\cup$, $\cap$, $\emptyset$ \\
Arithmetic: $\pm$, $\mp$, $\times$, $\div$, $\cdot$ \\
Logic: $\forall$, $\exists$, $\neg$, $\land$, $\lor$, $\Rightarrow$, $\Leftrightarrow$

\section{Aligned Equations}

Use \texttt{align*} to align multiple equations at \texttt{\&}:
\begin{align*}
    f(x)  &= x^2 + 2x + 1 \\
          &= (x + 1)^2
\end{align*}

\section{Matrices}

\begin{equation*}
    A = \begin{pmatrix} a & b \\ c & d \end{pmatrix}
    \qquad
    B = \begin{bmatrix} 1 & 0 \\ 0 & 1 \end{bmatrix}
\end{equation*}

\end{document}
